\documentclass[14pt,dvipdfmx]{jsarticle}
\input{lecture-preamble.tex}

\setcounter{definition}{3}
\setcounter{theorem}{8}
\setcounter{technique}{3}

\begin{document}
\lectureheader{数学Ⅱ}{5}{指数関数と対数関数}{2}{対数関数}{5-B}{常用対数の応用}{170--171}{常用対数を使って, 大きな数の桁数を求めよう}

% ============ 1ページ目（表・左） ============
\vspace{-1.6zw}
\heading{〜前時の復習〜}
\vspace{-0.35em}
\begin{reviewbox}{常用対数}
  \begin{enumerate}[label=\textbf{\arabic*}, leftmargin=1.5zw, itemsep=0.02em]
    \item $10$ を底とする対数 $\log_{10} M$ を\fitblankbf{\text{常用対数}}という.
    \item $M = a \times 10^{n}$ （$1 \leqq a < 10$）のとき $\log_{10} M = \fitblankbf{n + \log_{10} a}$.
    \item $0 \leqq \log_{10} a < \fitblankbf{1}$ である.
  \end{enumerate}
\end{reviewbox}

\vspace{0.2em}
\heading{常用対数の応用}
自然数 $N$ の桁数と, 常用対数 $\log_{10} N$ の関係を調べよう.

\vspace{0.15em}
\noindent{\bfseries 例\,13}\quad
自然数 $N$ が $3$ 桁の数のとき, $N$ のとり得る値の範囲を考える.
\par\vspace{0.12em}
$3$ 桁の自然数のうち, 最小の数は $100$, 最大の数は $999$ であるから
\par\vspace{0.08em}
{\centering
$100 \leqq N < 1000$\qquad すなわち\qquad
$10^{\fitblankbf{2}} \leqq N < 10^{\fitblankbf{3}}$
\par}
\vspace{0.08em}
両辺の常用対数をとると
\par\vspace{0.06em}
{\centering
$\log_{10} 10^{2} \leqq \log_{10} N < \log_{10} 10^{3}$\qquad よって\qquad
$\fitblankbf{2} \leqq \log_{10} N < \fitblankbf{3}$
\par}

\vspace{0.2em}
\begin{theorembox}{桁数と常用対数}
  自然数 $N$ が $k$ 桁の数であることは, 次と同値である.
  \par\vspace{0.12em}
  {\centering
  $10^{k-1} \leqq N < 10^{k}$\qquad
  $\Longleftrightarrow$\qquad
  $k-1 \leqq \log_{10} N < k$
  \par}
\end{theorembox}

\vspace{0.1em}
\textbf{注意}\quad $\log_{10} N$ の整数部分が $k-1$ のとき, $N$ は\fitblankbf{k}桁の数である.\\
すなわち, 桁数は $\log_{10} N$ の整数部分に $1$ を足した数である.

\vfill
\nextpage

% ============ 2ページ目（表・右） ============
\vspace{-1.4zw}
例13のように対数を用いると, $3^{20}$ などの大きな数の桁数を求められる.

\begin{exampleprob}{7}
  $3^{20}$ は何桁の数か。ただし, $\log_{10} 3 = 0.4771$ とする。
\end{exampleprob}

\begin{answerbox}{}[11]
  $\log_{10} 3^{20} = 20\log_{10} 3 = 20 \times 0.4771 = 9.5420$\\
  $9 < 9.5420 < 10$ であるから\\
  $10^{9} < 3^{20} < 10^{10}$
  \boxanswer{$10$ 桁}
\end{answerbox}

\vspace{0.12em}
\begin{techniquebox}{累乗の桁数の求め方}
  \begin{enumerate}[label=\textbf{\arabic*}, leftmargin=1.5zw, itemsep=0.02em]
    \item $\log_{10} a^{n} = \fitblankbf{n \log_{10} a}$ を計算する.
    \item $k-1 \leqq \log_{10} a^{n} < k$ となる整数 $k$ を求める.
    \item $a^{n}$ は\fitblankbf{k}桁の数である.
  \end{enumerate}
\end{techniquebox}

\vspace{0.15em}
\begin{practiceprob}{25}
  次の数は何桁か。ただし, $\log_{10} 2 = 0.3010$ とする。
  \begin{tasks}[label=(\arabic*), label-width=1.6zw, item-indent=2zw, after-item-skip=0.06em](2)
    \task $2^{20}$
    \task $2^{30}$
  \end{tasks}
\end{practiceprob}

\noindent
\begin{minipage}[t]{0.48\linewidth}
\begin{answerbox}{(1)}[8]
  $20\log_{10} 2 = 6.0200$\\
  $6 < 6.0200 < 7$
  \boxanswer{$7$ 桁}
\end{answerbox}
\end{minipage}\hfill
\begin{minipage}[t]{0.48\linewidth}
\begin{answerbox}{(2)}[8]
  $30\log_{10} 2 = 9.0300$\\
  $9 < 9.0300 < 10$
  \boxanswer{$10$ 桁}
\end{answerbox}
\end{minipage}

\vfill
\nextpage

% ============ 3ページ目（裏・左） ============
\vspace{-1.4zw}
自然数 $n$ に対して $2^{n}$ の桁数が与えられたとき,
常用対数を用いて $n$ の値を求めることができる.

\vspace{0.1em}
\noindent{\bfseries 応用例題\,3}\quad
$2^{n}$ が $10$ 桁の数となるような自然数 $n$ をすべて求めよ。
ただし, $\log_{10} 2 = 0.3010$ とする。

\vspace{0.08em}
\textbf{考え方}\quad
$2^{n}$ が $10$ 桁の数のとき, $10^{9} \leqq 2^{n} < 10^{10}$ が成り立つ.
常用対数をとると $n$ の $1$ 次不等式が得られる.

\begin{answerbox}{}[16]
  \makebox[13zw][l]{$2^{n}$ が $10$ 桁の数のとき}$10^{9} \leqq 2^{n} < 10^{10}$\\
  \makebox[13zw][l]{常用対数をとると}$9 \leqq n\log_{10} 2 < 10$\\
  \makebox[13zw][l]{$\log_{10} 2 = 0.3010 > 0$ であるから}$\dfrac{9}{0.3010} \leqq n < \dfrac{10}{0.3010}$ \quad ……①\\
  \makebox[13zw][l]{ここで}$9/0.3010=29.9\cdots$,\ $10/0.3010=33.2\cdots$\\
  \makebox[13zw][l]{①を満たす自然数 $n$ は}$n=30,\ 31,\ 32,\ 33$
\end{answerbox}

\vspace{0.12em}
\begin{techniquebox}{桁数が与えられたときの $n$ の求め方}
  \begin{enumerate}[label=\textbf{\arabic*}, leftmargin=1.5zw, itemsep=0.02em]
    \item $k$ 桁 $\Longleftrightarrow$ $10^{k-1} \leqq a^{n} < 10^{k}$ とおく.
    \item 常用対数をとって $n$ の不等式にする.
    \item $\log_{10} a > 0$ で両辺をわり, 整数 $n$ を求める.
  \end{enumerate}
\end{techniquebox}

\vfill
\nextpage

% ============ 4ページ目（裏・右） ============
\vspace{-1.2zw}
\begin{practiceprob}{26}
  $3^{n}$ が $8$ 桁の数となるような自然数 $n$ をすべて求めよ。
  ただし, $\log_{10} 3 = 0.4771$ とする。
\end{practiceprob}

\begin{answerbox}{}[14]
  \makebox[13zw][l]{$3^{n}$ が $8$ 桁の数のとき}$10^{7} \leqq 3^{n} < 10^{8}$\\
  \makebox[13zw][l]{常用対数をとると}$7 \leqq n\log_{10} 3 < 8$\\
  \makebox[13zw][l]{$\log_{10} 3 = 0.4771 > 0$ であるから}$\dfrac{7}{0.4771} \leqq n < \dfrac{8}{0.4771}$ \quad ……①\\
  \makebox[13zw][l]{ここで}$7/0.4771=14.6\cdots$,\ $8/0.4771=16.7\cdots$\\
  \makebox[13zw][l]{①を満たす自然数 $n$ は}$n=15,\ 16$
\end{answerbox}

\vspace{0.15em}
\textbf{注意}\quad 例題7と応用例題3は, どちらも次の関係を用いている.
\par\vspace{0.08em}
{\centering
$10^{k-1} \leqq N < 10^{k}$\quad $\Longleftrightarrow$\quad $N$ は $k$ 桁の数
\par}
\vspace{0.08em}
例題7は $\log_{10} N$ から桁数 $k$ を求め,
応用例題3は桁数 $k$ から $n$ を求める.

\vspace{0.25em}
\begin{summarybox}
  \begin{enumerate}[label=\textbf{\arabic*}, leftmargin=1.5zw, itemsep=0.06em]
    \item 自然数 $N$ が $k$ 桁 $\Longleftrightarrow$ $k-1 \leqq \log_{10} N < k$.
    \item 桁数は $\log_{10} N$ の整数部分に \fitblankbf{1} を足した数.
    \item $a^{n}$ の桁数は $\fitblankbf{n \log_{10} a}$ から求める.
  \end{enumerate}
\end{summarybox}

\vfill
\nexttime{常用対数を使って, 小数第何位かを調べてみよう!}

\end{document}
